\section{Related Work}

Retrieval-augmented models

Most models for fact-checking perform the following task: given a (statement, evidence) pair as input, they output a label indicating whether the statement is supported or refuted by the evidence (Thorne et al., 2018; Augenstein et al., 2019; Wadden et al., 2020). In many real-world situations, relevant evidence for a statement is not readily available, so some works have explored how to automatically retrieve evidence that may help support or refute a statement (Fan et al., 2020; Piktus et al., 2021).

More recent works have also explored how to “correct” statements so that they are more factually consistent with retrieved evidence (Shah et al., 2020; Thorne and Vlachos, 2021; Schuster et al., 2021; Logan IV et al., 2022; Schick et al., 2022). However, most of them used human-authored edits from FEVER (Thorne et al., 2018) as both their training and automatic evaluation data. FEVER’s claims were extracted from Wikipedia. This limits generalisability of these fact-checking models over generations of LLMs across various tasks and diverse domains.

Moreover, most previous work focus on sentence-level fact checking, which misaligns with long documents generated by recent LLMs: ChatGPT and GPT4, with multiple paragraphs. Our goal is to post-edit the output of any generation model across various domains over either long documents or short sentences. To this end, we will not be constrained by Wikipedia-specific resources for either training or evaluation, and we put emphasise on long documents.

RARR (Gao et al.,2022) is the most relevant work to ours. It handles the output of arbitrary TGMs in a unified framework by extending the research-and-revise paradigm, referring to as RARR. The goal to alleviate hallucinations by revising outputs of TGMs based on researched evidence (attribution report). However, the design of RARR framework still leads to several inherent limitations: (1) RARR would not be well-prepared to edit long documents; (2) RARR tends to preserve rather than delete claims that it cannot attribute. Some of these claims genuinely do not require attribution (e.g., opinions), but others are hallucination and should be removed; and (3) evaluation metric \textit{preservation: how much the revised text preserves aspects of the original text} is sometimes against our expectation of ideal revision.
For example, if the original output of a generation model is completely misguided or flawed, it can be difficult to revise the text without significant changes.

We also deal with text from different sources in a unified framework, either the output of arbitrary LLMs or misleading text written by human. While instead of retrieving evidence and editing over the whole document in RARR, we frame the whole procedure into seven subtasks shown in Figure 1. This pipeline naturally endows the ability to edit long documents: decomposing and decontextualising long document into independent atomic claims; and (2) flexible edits on subjective and fully-hallucinated statements: identifying check-worthiness of claims and flexible correction operations over claims, including complete deletion. 
It is necessary to develop new evaluation metrics and benchmarks for document-level fact checking.

We design the fact-checking framework as a plugin on top of rapidly-developed LLMs in the first stage of the project, without needing to re-design LLMs from scratch, and sought a built-in mechanism to address hallucinations in the second stage.



% Previous work handles outputs of different text generation models such as summarisation, question answering, dialog system and machine translation distinctively. 
% \citet{gao2022attributed} 
% it may serve as a natural baseline for future approaches.